% sections/05_other_regimes.tex -- Other regimes (lazy, rich non-incremental, stalled; discrete
% steps, SGD noise, finite samples). SKELETON.
% Sources: PLAN.md Q5 (lazy/rich parts) and Q8; briefs regime-map and optim-samples; results in
% experiments/regime_map/, experiments/optim_samples/ and their notes.
\section{Other regimes}\label{sec:other-regimes}

\placeholder{Skeleton. The regimes other than saddle-to-saddle, and how discrete steps, gradient
noise and finite samples modify the gradient-flow picture.}

\subsection{Lazy regime}\label{sec:other-lazy}
\placeholder{Large initialization in the centred and two-copy models: linearized dynamics, kernel
at initialization, time rescaling by $\isc^{2(\Kc-1)}$.}
% TODO: tangent kernel at init (theory_regimes.md); relative kernel change vs sigma; known in
% general (chizat2019lazy, jacot2018neural, woodworth2020kernel): mark \knownresult; the TTN kernel
% itself (its spectrum vs tree) may be new.

\subsection{Rich, non-incremental regime}\label{sec:other-rich}
\placeholder{Intermediate scales where features move but transitions overlap; what sets the
crossover.}
% TODO: overlapping events (rho_i <= 3); compare with DLN lazy-to-rich exact solutions (Domine et
% al.; kunin2024get; Tu, Aranguri & Jacot mixed regime).

\subsection{Stalled runs}\label{sec:other-stalled}
\placeholder{Runs that do not converge within the budget: long initial plateaus at small
$\isc$ and large $\Kc$, and the relation to barren plateaus.}
% TODO: escape time ~ sigma^{-(K-2)} makes small-init deep trees stall; contrast with barren
% plateau results for TN learning (liu2022presence; miao2024isometric; barthel2025absence).

\subsection{Large bond dimension: NTK versus \texorpdfstring{$\mu$P}{muP}}\label{sec:other-large-chi}
\placeholder{Per-core scalings for an infinite-bond-dimension limit with and without feature
learning.}
% TODO: spectral condition per core (Yang, Simon & Bernstein 2023; Yang & Hu 2021 abc
% parameterizations); express via per-core init scales and lr multipliers (Sec. 2.4 variant (e)).

\subsection{Discrete steps, gradient noise and finite samples}\label{sec:other-optim}
\placeholder{Gradient descent versus gradient flow, the stability limit and edge of stability;
SGD noise and plateau durations; sample complexity of the saddle-to-saddle and lazy regimes.}
% TODO(Q8, H6): GD stability limit eta_max ~ 2 / (K s_1^{2-2/K}) from the sharpness at the
% balanced minimum (verify or correct); EoS behaviour near eta_max (cohen2021gradient,
% lewkowycz2020large, even2023sgd); does a large step or SGD noise shorten plateaus at small sigma
% (plateau time vs batch size B and eta)? Finite n: test error vs n for small vs large init; low
% hierarchical rank bias of S2S vs min-norm tangent-space solution of lazy (razin2022implicit).
